\documentclass[a4paper,12pt]{article}
\usepackage[T2A]{fontenc}
\usepackage[utf8]{inputenc}
\usepackage[english,russian]{babel}
\usepackage{tikz}
\usetikzlibrary{shapes,arrows,positioning,calc}

\title{Лабораторная работа №8 \\ Рисование в LaTeX с помощью TikZ}
\author{Сунь Маосин}
\date{2026}

\begin{document}

\maketitle

\section{Цель работы}

Изучение возможностей пакета TikZ для создания векторной графики в LaTeX.

\section{Основные геометрические фигуры}

\subsection{Линии и стрелки}

\begin{tikzpicture}
  \draw (0,0) -- (2,0);
  \draw[->] (0,-0.5) -- (2,-0.5);
  \draw[<->] (0,-1) -- (2,-1);
  \draw[|-|] (0,-1.5) -- (2,-1.5);
\end{tikzpicture}

\subsection{Прямоугольники}

\begin{tikzpicture}
  \draw (0,0) rectangle (2,1);
  \draw[red, thick] (0,-1.5) rectangle (2,-0.5);
  \draw[blue, dashed] (0,-3) rectangle (2,-2);
  \fill[green] (3,0) rectangle (5,1);
  \filldraw[fill=yellow, draw=red] (3,-1.5) rectangle (5,-0.5);
\end{tikzpicture}

\subsection{Окружности и эллипсы}

\begin{tikzpicture}
  \draw (0,0) circle (0.5);
  \draw[blue] (2,0) circle (0.5 and 0.3);
  \fill[red] (0,-1.5) circle (0.3);
  \filldraw[fill=green, draw=black] (2,-1.5) circle (0.4);
\end{tikzpicture}

\section{Работа с цветом и стилями}

\begin{tikzpicture}
  \draw[red, thick] (0,0) -- (2,0);
  \draw[blue, very thick] (0,-0.5) -- (2,-0.5);
  \draw[green, dashed] (0,-1) -- (2,-1);
  \draw[orange, dotted] (0,-1.5) -- (2,-1.5);
  \draw[purple, dash dot] (0,-2) -- (2,-2);
\end{tikzpicture}

\section{Узлы (nodes)}

\subsection{Простые узлы}

\begin{tikzpicture}
  \node at (0,0) {Текст};
  \node[draw] at (2,0) {Рамка};
  \node[circle, draw] at (4,0) {Круг};
  \node[rectangle, draw, fill=yellow] at (6,0) {Жёлтый};
\end{tikzpicture}

\subsection{Узлы с соединениями}

\begin{tikzpicture}
  \node[draw] (A) at (0,0) {Узел A};
  \node[draw] (B) at (3,0) {Узел B};
  \draw[->] (A) -- (B);
  \draw[<->] (A) to[out=30,in=150] (B);
\end{tikzpicture}

\subsection{Цепочка узлов}

\begin{tikzpicture}
  \node[draw] (A) {Начало};
  \node[draw, right=of A] (B) {Шаг 1};
  \node[draw, right=of B] (C) {Шаг 2};
  \node[draw, below=of B] (D) {Ответвление};
  \draw[->] (A) -- (B);
  \draw[->] (B) -- (C);
  \draw[->] (B) -- (D);
\end{tikzpicture}

\section{Графики функций}

\subsection{Простой график}

\begin{tikzpicture}[scale=1.5]
  \draw[->] (-0.5,0) -- (3,0) node[right] {$x$};
  \draw[->] (0,-0.5) -- (0,3) node[above] {$y$};
  \draw[blue, thick] plot[domain=0:2.5] (\x, {0.5*\x*\x});
  \node at (1.5,2) {$y = \frac{1}{2}x^2$};
\end{tikzpicture}

\subsection{Тригонометрические функции}

\begin{tikzpicture}[scale=1.2]
  \draw[->] (-0.5,0) -- (7,0) node[right] {$x$};
  \draw[->] (0,-1.5) -- (0,1.5) node[above] {$y$};
  \draw[red, thick] plot[domain=0:6.28,smooth] (\x, {sin(\x r)});
  \draw[blue, thick] plot[domain=0:6.28,smooth] (\x, {cos(\x r)});
  \node[red] at (5,0.8) {$\sin x$};
  \node[blue] at (5,-0.8) {$\cos x$};
\end{tikzpicture}

\section{Циклы (foreach)}

\subsection{Повторяющиеся элементы}

\begin{tikzpicture}
  \foreach \x in {0,1,2,3,4}
    \draw (\x,0) circle (0.3);
  
  \foreach \x/\c in {0/red,1/green,2/blue,3/orange,4/purple}
    \fill[\c] (\x,-1) circle (0.3);
\end{tikzpicture}

\subsection{Координатная сетка}

\begin{tikzpicture}
  \draw[gray, very thin] (0,0) grid (4,3);
  \foreach \x in {0,1,2,3,4}
    \draw (\x,0.1) -- (\x,-0.1) node[below] {$\x$};
  \foreach \y in {0,1,2,3}
    \draw (0.1,\y) -- (-0.1,\y) node[left] {$\y$};
\end{tikzpicture}

\section{Сложные примеры}

\subsection{Простой граф}

\begin{tikzpicture}[
  roundnode/.style={circle, draw=blue!60, fill=blue!5, very thick},
  squarednode/.style={rectangle, draw=red!60, fill=red!5, very thick}
]
  \node[roundnode] (a) {A};
  \node[squarednode] (b) [right=of a] {B};
  \node[roundnode] (c) [below right=of b] {C};
  
  \draw[->] (a) -- (b);
  \draw[->] (b) -- (c);
  \draw[->, bend left] (a) to (c);
\end{tikzpicture}

\subsection{Схема алгоритма}

\begin{tikzpicture}[
  startstop/.style={rectangle, draw, fill=red!20, text width=2cm, align=center},
  io/.style={trapezium, trapezium left angle=70, trapezium right angle=110, draw, fill=blue!20, text width=2cm, align=center},
  process/.style={rectangle, draw, fill=orange!20, text width=2.5cm, align=center},
  decision/.style={diamond, draw, fill=green!20, text width=2cm, align=center},
  arrow/.style={->, thick}
]

  \node[startstop] (start) {Начало};
  \node[io, below=of start] (input) {Ввод данных};
  \node[process, below=of input] (process1) {Обработка};
  \node[decision, below=of process1] (decide) {Условие};
  \node[process, right=of decide] (process2) {Действие 2};
  \node[startstop, below=of decide] (stop) {Конец};

  \draw[arrow] (start) -- (input);
  \draw[arrow] (input) -- (process1);
  \draw[arrow] (process1) -- (decide);
  \draw[arrow] (decide) -- node[anchor=east] {Да} (stop);
  \draw[arrow] (decide) -- node[anchor=south] {Нет} (process2);
  \draw[arrow] (process2) |- (process1);
\end{tikzpicture}

\section{Вывод}

В ходе работы были изучены:
\begin{itemize}
  \item основные геометрические фигуры в TikZ;
  \item работа с цветами и стилями линий;
  \item создание и соединение узлов;
  \item построение графиков функций;
  \item использование циклов для повторяющихся элементов;
  \item создание сложных схем и графов.
\end{itemize}

\end{document}